\documentclass{article}

\usepackage[utf8]{inputenc}     % pdfLaTeX
\usepackage[T1]{fontenc}
\usepackage[magyar]{babel}
\usepackage{amssymb}
\usepackage{amsmath}

\title{Magyar Egyházi Irodalom}
\author{Lengyel Zoltán Péter, Rovenszky Robin, Simon László}
\date{Legutóbb frissítve: 2026. Június 02.}

\begin{document}

\maketitle
\tableofcontents
\newpage

\section{A Biblia}
insert Bible stuff here 


%LACI go take notes
%nah
%bro you are uselessssssssssssss
%didnt ask + dont care
\section{Prédikáció magyar nyelven}
\begin{itemize}
    \item egyházi beszéd \begin{itemize}
        \item Jézus életéről szóló tanítás \begin{itemize}
            \item PRÓZAI IRODALOM
        \end{itemize}
    \end{itemize}
    \item Egyszerű, köznapi történeten keresztül tanít
\end{itemize}

\subsection{Halotti beszéd és könyörgés}
\begin{itemize}
    \item Első fennmaradt magyar nyelvű szövegemlék - 1195 körül
    \item Pray-kódex-ben találták meg
    \item temetési beszéd
    \item két részből áll \begin{itemize}
        \item ima - forma kötöttebb
        \item beszéd - E/2, közvetlenebb
    \end{itemize}
    \item por és hamu 
    \item "Halálnak halálával halsz" - nyomatékosít, szótőismétlés/alliteráció figura etimológika
    
\end{itemize}

\section{Ómagyar Mária siralom}
\begin{itemize}
    \item Mária felajánlja életét fia életéért cserébe
    \item Szerzőség kérdése
    \item Leuveny kódex
    \item siralom mint műfaj
    \item drámai jelenetezés
    \item virág szimbolika 
\end{itemize}

\section{Lovagi költészet}
\begin{itemize}
    \item XII-XIII, század Nyugat-Európa
    \item keresztes hadjáratokkal jelennek meg a lovagok
    \item társadalmi kódrendszer - lovag fogalma \begin{itemize}
        \item norma, értékrendszer \\ $\implies$ bátor, vakmerő, intelligens, művelt, eszes \\ + morális értékek $\implies$ igazságosság, becsületesség, őszinteség
    \end{itemize}
    \item fogadalmak \begin{itemize}
        \item nők, elesettek, gyámoltalanok védelme
        \item hűség az urához, országhoz, haza védelme!
        \item isten szolgálata
    \end{itemize}
    \item Lovagi irodalom \begin{itemize}
        \item világi irodalom ($\neq$ egyház) \\ $\implies$ témái: \\
        \begin{tabular}{c c c}
             égi szerelem & vs & földi szerelem  \\
             elérhetetlen, idilli & & elérhető \\ harmonikus, álomszerű & & testi vágyak \\lelki szépség & & \\
             \\
             trubadúr & & Minnesänger \\
            (Fr. o.) & & (Németalföld)
             
        \end{tabular}
    \end{itemize}
\end{itemize}
\end{document}